\documentclass[11pt,a4paper]{article}

\usepackage{iftex}
\ifXeTeX
  \usepackage{fontspec}
  \setmainfont{Times New Roman}
\else\ifLuaTeX
  \usepackage{fontspec}
  \setmainfont{Times New Roman}
\else
  \usepackage[utf8]{inputenc}
  \usepackage[vietnamese]{babel}
  \usepackage[T5,T1]{fontenc}
  \usepackage{mathptmx}
\fi\fi

\usepackage{tabularx}
\usepackage{booktabs}
\usepackage{array}
\newcolumntype{L}[1]{>{\raggedright\arraybackslash}p{#1}}
\newcolumntype{C}[1]{>{\centering\arraybackslash}p{#1}}
\newcolumntype{R}[1]{>{\raggedleft\arraybackslash}p{#1}}
\newcolumntype{Y}{>{\raggedright\arraybackslash}X}
\newcolumntype{Z}{>{\centering\arraybackslash}X}
\usepackage[top=2.2cm,bottom=2.2cm,left=2.2cm,right=2.2cm]{geometry}
\usepackage{amsmath,amssymb}
\usepackage{longtable}
\usepackage{multirow}
\usepackage{fancyhdr}
\usepackage{titlesec}
\usepackage{float}
\usepackage{enumitem}
\usepackage{tcolorbox}
\usepackage{tikz}
\usetikzlibrary{shapes.geometric, arrows, positioning}

% Hyperref hoàn toàn đen trắng, không viền, không màu
\usepackage[hidelinks]{hyperref}

% Cấu hình Header và Footer đen trắng
\setlength{\headheight}{14pt}
\pagestyle{fancy}
\fancyhf{}
\fancyhead[L]{\small\textit{Dự án LMS Tiên tiến \& Learning Analytics}}
\fancyhead[R]{\small\textit{Báo cáo tuần 2}}
\fancyfoot[C]{\thepage}
\renewcommand{\headrulewidth}{0.5pt}

% Định dạng tiêu đề hoàn toàn chữ đen chuẩn mực
\titleformat{\section}
  {\normalfont\Large\bfseries}{\thesection.}{0.8em}{}
\titleformat{\subsection}
  {\normalfont\large\bfseries}{\thesubsection.}{0.6em}{}
\titleformat{\subsubsection}
  {\normalfont\normalsize\bfseries}{\thesubsubsection.}{0.5em}{}

% Hộp ghi chú / đóng khung thuần đen trắng
\newtcolorbox{monobox}[1][]{
  colback=white,
  colframe=black,
  fonttitle=\bfseries,
  coltitle=black,
  colbacktitle=white,
  title=#1,
  boxrule=0.8pt,
  arc=0mm,
  left=3mm, right=3mm, top=2.5mm, bottom=2.5mm
}

\begin{document}

% Trang bìa / Tiêu đề học thuật đen trắng
\begin{center}
    
    {\Large \textbf{BÁO CÁO NGHIÊN CỨU TUẦN 2}}\\[0.4cm]
    {\LARGE \textbf{LMS TIÊN TIẾN VÀ LEARNING ANALYTICS:\\[0.2cm] PHÂN TÍCH DỮ LIỆU HỌC TẬP VÀ DỰ ĐOÁN NGUY CƠ}}\\[0.6cm]
    
    {\normalsize \textbf{Mục tiêu}: Thu thập, nạp và kiểm định toàn diện chất lượng bộ dữ liệu OULAD}\\[1cm]
    
    \begin{tabular}{rl}
        \textbf{Sinh viên thực hiện:} & Vũ Hoà Bình \\
        \textbf{Mã số sinh viên:} & 202416137 \\
        \textbf{Bộ dữ liệu nghiên cứu:} & Open University Learning Analytics Dataset (OULAD) \\
        \textbf{Thời gian hoàn thành:} & 22 / 10 / 2026 \\
    \end{tabular}
\end{center}

\vspace{0.8cm}
\hrule
\vspace{0.5cm}

\tableofcontents
\vspace{0.5cm}
\hrule
\vspace{0.8cm}

\newpage

\section{Tổng quan tiến độ tuần 2}

Trong tuần 2, trọng tâm chuyển từ cơ sở lý thuyết sang thực nghiệm dữ liệu thực tế.

\begin{table}[H]
\centering
\small
\renewcommand{\arraystretch}{1.3}
\begin{tabularx}{\textwidth}{@{}>{\centering\arraybackslash}p{1.2cm} X >{\centering\arraybackslash}p{3cm}@{}}
\toprule
\textbf{STT} & \textbf{Nhiệm vụ nghiên cứu \& triển khai} & \textbf{Trạng thái} \\
\midrule
\textbf{2.1} & Tải và giải nén hoàn chỉnh 7 file CSV OULAD vào thư mục \texttt{data/raw/} & HOÀN THÀNH \\
\textbf{2.2} & Xây dựng pipeline nạp dữ liệu bằng Pandas, kiểm tra cấu trúc \& kiểu dữ liệu & HOÀN THÀNH \\
\textbf{2.3} & Kiểm định giá trị khuyết thiếu (Missing Values Audit) và phân tích nguyên nhân & HOÀN THÀNH \\
\textbf{2.4} & Kiểm tra tính toàn vẹn quan hệ (Referential Integrity) giữa 7 bảng liên kết & HOÀN THÀNH \\
\textbf{2.5} & Khảo sát phân bố biến mục tiêu (\texttt{final\_result}) và ánh xạ bài toán nhị phân & HOÀN THÀNH \\
\textbf{2.6} & Khảo sát ban đầu tương tác tài nguyên LMS theo Khung lý thuyết SEDL & HOÀN THÀNH \\
\bottomrule
\end{tabularx}
\caption{Bảng tổng kết tiến độ hoàn thành các nhiệm vụ trong Tuần 2}
\end{table}

\section{Kiến trúc vật lý \& thống kê kỹ thuật 7 bảng OULAD}


Dưới đây là bảng tổng hợp các thông số kỹ thuật, số lượng bản ghi, số chiều đặc trưng và mức tiêu thụ bộ nhớ RAM khi nạp vào môi trường phân tích:

\begin{table}[H]
\centering
\small
\renewcommand{\arraystretch}{1.25}
\begin{tabularx}{\textwidth}{@{}l c c r r X@{}}
\toprule
\textbf{Bảng dữ liệu} & \textbf{Số dòng} & \textbf{Số cột} & \textbf{Dung lượng đĩa} & \textbf{RAM nạp} & \textbf{Nội dung nghiệp vụ} \\
\midrule
\texttt{courses.csv} & 22 & 3 & 0.001 MB & 0.00 MB & 22 khóa học (7 modules $\times$ 4 kỳ) \\
\texttt{assessments.csv} & 206 & 6 & 0.010 MB & 0.01 MB & Cấu trúc bài kiểm tra (TMA, CMA, Exam) \\
\texttt{studentInfo.csv} & 32,593 & 12 & 3.300 MB & 4.85 MB & Nhân khẩu học \& kết quả cuối kỳ \\
\texttt{studentRegistration.csv} & 32,593 & 5 & 1.060 MB & 1.49 MB & Nhật ký ngày đăng ký và hủy môn \\
\texttt{studentAssessment.csv} & 173,912 & 5 & 5.430 MB & 6.63 MB & Điểm số và ngày nộp bài thực tế \\
\texttt{vle.csv} & 6,364 & 6 & 0.250 MB & 0.38 MB & Danh mục tài nguyên Moodle \\
\texttt{studentVle.csv} & 10,655,280 & 6 & 432.810 MB & 569.05 MB & Log tương tác click chuột theo ngày \\
\midrule
\textbf{Tổng cộng} & \textbf{10,868,470} & \textbf{43} & \textbf{442.861 MB} & \textbf{582.41 MB} & \textbf{Hệ thống 7 bảng liên kết} \\
\bottomrule
\end{tabularx}
\caption{Thống kê kỹ thuật và mức tiêu thụ tài nguyên của 7 bảng dữ liệu OULAD}
\end{table}

\newpage

\section{Kiểm định chất lượng dữ liệu \& Missing Data Audit}

\subsection{Kết quả kiểm tra toàn diện giá trị khuyết thiếu}

\begin{table}[H]
\centering
\small
\renewcommand{\arraystretch}{1.25}
\setlength{\tabcolsep}{4.5pt}
\begin{tabularx}{\textwidth}{@{}l r r r r Y@{}}
\toprule
\textbf{Bảng dữ liệu} & \textbf{Tổng số ô} & \textbf{Số ô Missing} & \textbf{Tỷ lệ (\%)} & \textbf{Dòng trùng} & \textbf{Đặc điểm dữ liệu} \\
\midrule
\texttt{courses} & 66 & 0 & 0.00\% & 0 & Hoàn hảo, không khuyết thiếu \\
\texttt{assessments} & 1,236 & 11 & 0.89\% & 0 & Thiếu ngày thi Exam cuối kỳ \\
\texttt{studentInfo} & 391,116 & 1,111 & 0.28\% & 0 & Thiếu chỉ số kinh tế vùng IMD \\
\texttt{studentRegistration} & 162,965 & 22,566 & 13.85\% & 0 & Missing có ý nghĩa cấu trúc \\
\texttt{studentAssessment} & 869,560 & 173 & 0.02\% & 0 & Thiếu điểm nộp bài của sinh viên \\
\texttt{vle} & 38,184 & 10,486 & 27.46\% & 0 & Tài nguyên mở không giới hạn tuần \\
\texttt{studentVle} & 63,931,680 & 0 & 0.00\% & 787,170* & Các phiên click trong ngày \\
\bottomrule
\end{tabularx}
\caption{Bảng tổng hợp kiểm định giá trị khuyết thiếu (Missing Values Audit)}
\end{table}

Trong bảng \texttt{studentVle}, 787,170 bản ghi trùng khóa \texttt{(id\_student, id\_site, date)} phản ánh các phiên tương tác riêng biệt trong ngày. Khi tổng hợp đặc trưng cho mô hình, các dòng này sẽ được cộng gộp tự nhiên theo $\sum \texttt{sum\_click}$.

\subsection{Phân tích bản chất và chiến lược xử lý dữ liệu khuyết thiếu}

Để đảm bảo tính hợp lệ sư phạm và độ tin cậy của mô hình học máy, các biến khuyết thiếu được phân loại thành 4 nhóm bản chất:

\begin{enumerate}[leftmargin=*, itemsep=2pt, topsep=2pt]
    \item \textbf{Khuyết thiếu có ý nghĩa cấu trúc nghiệp vụ (Structural Missing) -- Cảnh báo rò rỉ dữ liệu}:
    \begin{itemize}
        \item Cột \texttt{studentRegistration.date\_unregistration} có \textbf{22,521 giá trị khuyết thiếu (69.1\%)}.
        \item \textit{Bản chất}: Sinh viên không hủy môn học sẽ không có ngày hủy. Các sinh viên có ngày hủy môn đều có nhãn \texttt{final\_result == 'Withdrawn'}.
        \item\textit{Chiến lược xử lý}: \textbf{Không sử dụng cột này làm đặc trưng dự đoán} trong bất kỳ mốc thời gian nào để triệt tiêu hoàn toàn rủi ro rò rỉ thông tin tương lai (Data Leakage).
    \end{itemize}

    \item \textbf{Khuyết thiếu do đặc thù lịch trình thi cử}:
    \begin{itemize}
        \item Cột \texttt{assessments.date} thiếu \textbf{11 giá trị (5.3\%)}.
        \item \textit{Bản chất}: Toàn bộ 11 bài bị thiếu đều là bài thi kết thúc môn (\texttt{assessment\_type == 'Exam'}). Tại The Open University, ngày thi cuối kỳ được hội đồng đào tạo ấn định riêng trong tuần thi tập trung thay vì cố định từ đầu khóa.
        \item \textit{Chiến lược xử lý}: Gán hạn thi bằng độ dài khóa học \texttt{module\_presentation\_length} lấy từ bảng \texttt{courses}.
    \end{itemize}

    \item \textbf{Khuyết thiếu do chính sách mở của hệ thống LMS}:
    \begin{itemize}
        \item Hai cột \texttt{vle.week\_from} và \texttt{vle.week\_to} thiếu \textbf{5,243 giá trị (82.4\%)}.
        \item \textit{Bản chất}: Các tài nguyên học tập (như tài liệu \texttt{oucontent}, diễn đàn \texttt{forumng}, trang chủ \texttt{homepage}) được mở tự do xuyên suốt khóa học cho sinh viên tra cứu, không giới hạn tuần bắt đầu/kết thúc.
        \item \textit{Chiến lược xử lý}: Sử dụng loại tài nguyên (\texttt{activity\_type}) để phân nhóm đặc trưng SEDL.
    \end{itemize}

    \pagebreak
    \item \textbf{Khuyết thiếu thông tin cá nhân và điểm số}:
    \begin{itemize}
        \item Cột \texttt{studentInfo.imd\_band} thiếu \textbf{1,111 giá trị (3.4\%)}: Do sinh viên không cư trú tại vùng điều tra chỉ số nghèo đói của Vương quốc Anh hoặc không kê khai mã bưu chính. Xử lý bằng cách tạo nhóm danh mục riêng: \texttt{'Missing'}.
        \item Cột \texttt{studentAssessment.score} thiếu \textbf{173 giá trị (0.1\%)}: Do sinh viên nộp bài nhưng gặp sự cố kỹ thuật hoặc bài thi chưa hoàn tất quy trình chấm. Xử lý bằng cách gán giá trị 0 hoặc median theo bài thi.
    \end{itemize}
\end{enumerate}

\section{Kiểm tra tính toàn vẹn quan hệ}

\subsection{Kiểm tra liên kết khóa ngoại giữa các bảng}

Một hệ thống phân tích học tập đòi hỏi dữ liệu phải có tính nhất quán quan hệ tuyệt đối. Tiến hành kiểm tra đối chiếu tập hợp khóa ngoại:
\begin{itemize}[leftmargin=*, itemsep=2pt, topsep=2pt]
    \item \textbf{Kiểm tra 1}: Sinh viên trong bảng log tương tác \texttt{studentVle} không tồn tại trong bảng thông tin trung tâm \texttt{studentInfo}: $\mathbf{0}$ \textbf{sinh viên (100\% hợp lệ)}.
    \item \textbf{Kiểm tra 2}: Sinh viên có bài nộp trong \texttt{studentAssessment} không tồn tại trong \texttt{studentInfo}: $\mathbf{0}$ \textbf{sinh viên (100\% hợp lệ)}.
    \item \textbf{Kiểm tra 3}: Mã bài kiểm tra trong bảng \texttt{studentAssessment} không nằm trong danh mục bài thi của \texttt{assessments}: $\mathbf{0}$ \textbf{bài (100\% hợp lệ)}.
    \item \textbf{Kiểm tra 4}: Mã tài nguyên \texttt{id\_site} trong log tương tác không nằm trong danh mục tài nguyên \texttt{vle}: $\mathbf{0}$ \textbf{tài nguyên (100\% hợp lệ)}.
\end{itemize}


Dữ liệu OULAD đạt độ toàn vẹn quan hệ hoàn hảo (Referential Integrity = 100\%). Toàn bộ các bảng liên kết chặt chẽ qua các khóa chính--khóa ngoại, không phát sinh bất kỳ bản ghi mồ côi (orphan records) nào. Đây là tiền đề vững chắc để thực hiện các phép gom nhóm và nối dữ liệu (Merge/Join).


\subsection{Sơ đồ luồng nạp và kiểm định dữ liệu (Audit Pipeline)}

\begin{center}
\begin{tikzpicture}[node distance=1.1cm, auto,
    block/.style={rectangle, draw=black, fill=white, inner sep=4.5pt, text width=2.9cm, font=\scriptsize, align=center},
    decision/.style={diamond, draw=black, fill=white, inner sep=1.5pt, text width=1.9cm, font=\tiny, align=center},
    edge/.style={draw=black, thick, -latex}]

    \node[block] (raw) {\textbf{Kho dữ liệu thô}\\\tiny 7 tệp CSV (442.5 MB)\\\texttt{data/raw/}};
    \node[block, right=1.0cm of raw] (load) {\textbf{Nạp dữ liệu}\\\tiny \texttt{pandas.read\_csv}};
    \node[decision, right=1.0cm of load] (check) {\textbf{Kiểm tra Khóa}\\Tính toàn vẹn \& Missing};
    \node[block, above right=0.2cm and 1.0cm of check] (clean) {\textbf{Dữ liệu Hợp chuẩn}\\\tiny 0 bản ghi mồ côi\\Missing $< 0.3\%$};
    \node[block, below right=0.2cm and 1.0cm of check] (log) {\textbf{Cảnh báo Leakage}\\\tiny Cách ly \texttt{date\_unreg}\\Gán nhãn Missing IMD};

    \path[edge] (raw) -- (load);
    \path[edge] (load) -- (check);
    \path[edge] (check) |- node[above, font=\tiny]{Khóa ngoại khớp 100\%} (clean);
    \path[edge] (check) |- node[below, font=\tiny]{Structural Missing} (log);
\end{tikzpicture}
\end{center}

\newpage

\section{Phân tích phân bố biến mục tiêu (Target Distribution)}

\subsection{Phân bố 4 lớp kết quả học tập gốc}

Biến mục tiêu \texttt{final\_result} trong bảng \texttt{studentInfo} phản ánh kết quả học tập cuối cùng của 32,593 lượt đăng ký môn học:

\begin{table}[H]
\centering
\small
\renewcommand{\arraystretch}{1.25}
\begin{tabularx}{\textwidth}{@{}l c c X@{}}
\toprule
\textbf{Kết quả gốc (\texttt{final\_result})} & \textbf{Số lượng sinh viên} & \textbf{Tỷ trọng (\%)} & \textbf{Ý nghĩa sư phạm trong giáo dục từ xa} \\
\midrule
\textbf{Pass} (Đỗ môn học) & 12,361 & 37.9\% & Đạt yêu cầu tất cả bài kiểm tra và điểm thi \\
\textbf{Withdrawn} (Bỏ học giữa chừng) & 10,156 & 31.2\% & Rút môn học trước hoặc trong quá trình học \\
\textbf{Fail} (Trượt môn học) & 7,052 & 21.6\% & Hoàn thành khóa học nhưng điểm không đạt chuẩn \\
\textbf{Distinction} (Đỗ xuất sắc) & 3,024 & 9.3\% & Đạt điểm số xuất sắc ở toàn bộ các cấu phần \\
\midrule
\textbf{Tổng số} & \textbf{32,593} & \textbf{100.0\%} & \textbf{Quy mô toàn bộ bộ dữ liệu OULAD} \\
\bottomrule
\end{tabularx}
\caption{Phân bố chi tiết 4 lớp kết quả học tập gốc trong OULAD}
\end{table}

\subsection{Ánh xạ và phân tích bài toán phân loại nhị phân}

Nhằm phục vụ bài toán can thiệp sớm của cố vấn học tập, 4 lớp gốc được ánh xạ thành bài toán nhị phân chuẩn:
\begin{itemize}
    \item \textbf{Lớp 0 -- Not At-Risk (An toàn)}: Gồm \texttt{Pass} và \texttt{Distinction} với tổng số \textbf{15,385 sinh viên (chiếm 47.2\%)}.
    \item \textbf{Lớp 1 -- At-Risk (Nguy cơ)}: Gồm \texttt{Fail} và \texttt{Withdrawn} với tổng số \textbf{17,208 sinh viên (chiếm 52.8\%)}.
\end{itemize}


\subsection{Sự phân hóa hành vi cốt lõi giữa nhóm Withdrawn và Fail}

Mặc dù cùng được gán nhãn \textbf{At-Risk}, phân tích khám phá chỉ ra sự khác biệt sâu sắc giữa hai phân nhóm này:
\begin{enumerate}
    \item \textbf{Nhóm Withdrawn (Bỏ học giữa kỳ - 31.2\%)}: Có xu hướng ngừng click đột ngột trên hệ thống LMS ngay sau 20\%--30\% thời lượng môn học. Đặc trưng nổi bật là khoảng cách ngưng hoạt động (\texttt{inactivity\_gap}) tăng vọt.
    \item \textbf{Nhóm Fail (Thi trượt - 21.6\%)}: Vẫn duy trì đăng nhập và click chuột đều đặn đến cuối kỳ nhưng điểm số bài tập (\texttt{score}) thấp và thời gian nộp bài bị trễ hạn thường xuyên (\texttt{submit\_delay}).
\end{enumerate}
Phát hiện này củng cố tính đúng đắn của việc kết hợp đa chiều các đặc trưng SEDL (Behavioural kết hợp Cognitive) thay vì chỉ nhìn vào số lượng click đơn thuần.

\newpage

\section{Phân tích tương tác tài nguyên học tập theo Khung SEDL}

\subsection{Tỷ trọng tương tác các loại tài nguyên VLE}
Bảng log \texttt{studentVle} ghi nhận tổng cộng \textbf{10,655,280 lượt click} trên 20 danh mục tài nguyên trực tuyến. Thống kê tỷ lệ phân bổ click cho thấy sự tập trung rõ rệt vào 4 nhóm tài nguyên chủ lực:

\begin{table}[H]
\centering
\small
\renewcommand{\arraystretch}{1.25}
\begin{tabularx}{\textwidth}{@{}l c c X@{}}
\toprule
\textbf{Loại tài nguyên (\texttt{activity\_type})} & \textbf{Tỷ lệ lượt click} & \textbf{Chiều tương tác SEDL} & \textbf{Vai trò sư phạm trong hệ thống VLE} \\
\midrule
\texttt{oucontent} & $\approx 50.2\%$ & Cognitive & Tài liệu bài giảng số hóa chính thống của môn học \\
\texttt{forumng} & $\approx 24.8\%$ & Social & Diễn đàn trao đổi thảo luận giữa SV và giảng viên \\
\texttt{homepage} & $\approx 12.1\%$ & Behavioural & Trang chủ môn học, kiểm tra thông báo và lịch trình \\
\texttt{quiz} & $\approx 8.3\%$ & Cognitive & Bài tập trắc nghiệm tự luyện kiểm tra kiến thức \\
\texttt{subpage}, \texttt{resource}, \texttt{url} & $\approx 4.6\%$ & Cognitive/Support & Tài liệu đọc thêm, liên kết ngoài và tệp đính kèm \\
\bottomrule
\end{tabularx}
\caption{Phân bố lượt click theo loại tài nguyên VLE và ánh xạ theo Khung SEDL}
\end{table}

\subsection{Ý nghĩa đối với quá trình trích xuất đặc trưng (Feature Engineering)}
Kết quả phân tích tương tác tài nguyên khẳng định tính tương thích cao của OULAD với Khung lý thuyết SEDL (Bergdahl et al., 2024):
\begin{itemize}
    \item Tài nguyên học tập chiếm hơn \textbf{50\%} tổng số click, chứng tỏ sinh viên dành phần lớn nỗ lực nhận thức (\textit{Cognitive Engagement}) để nghiên cứu học liệu.
    \item Diễn đàn trao đổi chiếm gần \textbf{25\%} tổng số tương tác, là bằng chứng thực nghiệm rõ ràng để đo lường mức độ hòa nhập cộng đồng (\textit{Social Engagement}).
    \item Tỷ lệ làm quiz tự luyện ($\sim 8\%$) sẽ là đặc trưng phản ánh năng lực tự điều chỉnh học tập (\textit{self-regulated learning}).
\end{itemize}

\newpage

\section{Kế hoạch triển khai tuần 3}


Bước sang tuần 3, đề tài sẽ tiến hành khai phá dữ liệu học tập (Exploratory Data Analysis - EDA) để trả lời định lượng cho các Giả thuyết H1a, H1b, H1c trong Câu hỏi nghiên cứu RQ1:

\begin{table}[H]
\centering
\small
\renewcommand{\arraystretch}{1.25}
\begin{tabularx}{\textwidth}{@{}>{\centering\arraybackslash}p{1.2cm} X@{}}
\toprule
\textbf{STT} & \textbf{Nội dung công việc tuần 3} \\
\midrule
\textbf{3.1} & Phân tích phân bố số lượt click giữa nhóm At-Risk vs Not At-Risk \\
\textbf{3.2} & Khảo sát chuỗi thời gian: Xu hướng click theo tuần học (Temporal Trends) \\
\textbf{3.3} & Đánh giá mối quan hệ giữa điểm bài tập đầu tiên (\texttt{TMA1}) với kết quả môn \\
\textbf{3.4} & Kiểm định thống kê: T-test và ANOVA cho sự khác biệt mức độ tương tác \\
\textbf{3.5} & Phân tích ma trận tương quan giữa các nhóm đặc trưng nhân khẩu học và LMS \\
\bottomrule
\end{tabularx}
\caption{Kế hoạch phân công nhiệm vụ nghiên cứu trong tuần 3}
\end{table}

\vspace{0.8cm}
\begin{center}
    \textbf{--- HẾT ---}
\end{center}

\end{document}
